\chapter{Complessità computazionale}
\label{chap:complessita}

L'esistenza di un algoritmo risolutivo, discussa nella
\cref{sec:fi-algoritmi}, non ne determina l'efficienza. La
\emph{complessità computazionale} studia le risorse necessarie
all'esecuzione in funzione della dimensione dell'ingresso: principalmente
tempo e memoria. L'analisi permette di confrontare algoritmi che
risolvono lo stesso problema e di prevederne il comportamento per
ingressi non ancora elaborati.\footnote{Materiale di riferimento:
S. Canazza, \emph{Complessità computazionale}, dispense del corso di
Fondamenti di Informatica.}

\section{Analisi sperimentale}
\label{sec:fi-complessita-sperimentale}

L'analisi sperimentale misura le prestazioni di un'implementazione
eseguendola su ingressi di dimensione diversa. Se il tempo impiegato per
un ingresso di dimensione \(n\) è \(t(n)\), il rapporto fra \(t(2n)\) e
\(t(n)\) descrive l'effetto del raddoppio della dimensione.

Nel modello \(t(n)=an^b\), con \(a>0\), si ha
\[
  \frac{t(2n)}{t(n)}=2^b,
  \qquad
  b=\log_2\!\left(\frac{t(2n)}{t(n)}\right).
\]
Un rapporto prossimo a \(4\) suggerisce quindi una crescita quadratica;
un rapporto prossimo a \(2\) suggerisce una crescita lineare. La
stima non costituisce, da sola, una dimostrazione della complessità.

In Python, il modulo \texttt{time} offre \texttt{time.time()}, che
restituisce un istante espresso in secondi, e
\texttt{time.perf\_counter()}, adatto alla misura di intervalli di
tempo. La durata si ottiene sottraendo le letture effettuate prima e
dopo l'esecuzione. Per confrontare più dimensioni si possono registrare
i risultati in formato CSV, utilizzando i file di testo introdotti
nella \cref{sec:fi-python-file}. Un esempio completo di misurazione
sarà presentato dopo le due implementazioni di \texttt{reverse}.

Il tempo misurato dipende anche dai valori dei dati, dall'hardware,
dall'interprete o compilatore, dal linguaggio, dall'implementazione e
dalle condizioni di esecuzione. Ripetere gli esperimenti riduce
l'influenza delle fluttuazioni, ma non elimina tali dipendenze. Per
confrontare gli algoritmi indipendentemente da una particolare macchina
si ricorre all'analisi teorica.

\section{Analisi teorica e modello di costo}
\label{sec:fi-complessita-modello}

L'analisi teorica associa un costo alle operazioni eseguite, senza
misurarne direttamente la durata. Nel \emph{modello di costo unitario},
le operazioni primitive hanno costo costante: assegnamenti,
operazioni aritmetiche e logiche, confronti e accessi a una posizione
di una lista. Il costo dell'accesso per indice è giustificato
dall'organizzazione mediante array della \cref{sec:fi-strutture-array}.
Si tratta di un'astrazione, che presuppone operandi di
dimensione limitata e operazioni elementari eseguibili in tempo
costante.

Non tutte le istruzioni Python hanno costo unitario. Su una lista,
\texttt{len(lista)} richiede tempo costante, mentre la costruzione di
\texttt{[0] * n} richiede tempo proporzionale a \(n\). Anche il costo
di una funzione dipende dalle operazioni svolte al suo interno e non
può essere identificato con quello della sola chiamata.

Fissato il modello, si indica con \(T(x)\) il costo di esecuzione
sull'ingresso \(x\). Per ingressi di dimensione \(n\), il costo può
dipendere anche dal contenuto: si possono pertanto considerare il
caso peggiore e il caso migliore. Nei programmi esaminati di seguito
il numero delle operazioni considerate dipende soltanto dalla
lunghezza \(n\) della lista.

Nel conteggio dei passi si sommano i costi dei blocchi eseguiti in
sequenza. Per un'istruzione condizionale si considera il ramo
effettivamente percorso; per un limite superiore si può maggiorare
il costo con quello del ramo più oneroso. Il costo di un ciclo è la
somma dei costi delle sue iterazioni, che possono essere differenti.
Il modello può essere semplificato ulteriormente scegliendo come passi
base soltanto le operazioni rilevanti per il confronto in esame.

La \emph{complessità temporale} riguarda il numero di operazioni;
quella \emph{spaziale} riguarda la quantità di memoria richiesta. Le
due risorse devono essere valutate separatamente.

\section{Notazioni asintotiche}
\label{sec:fi-complessita-asintotica}

Per confrontare le funzioni di costo si studia il loro andamento per
\(n\) sufficientemente grande. Le notazioni asintotiche trascurano i
fattori moltiplicativi costanti e il comportamento su un numero finito
di dimensioni iniziali.

\begin{definition}[Notazioni \(O\), \(\Omega\) e \(\Theta\)]
\label{def:fi-notazioni-asintotiche}
Siano \(f,g\colon\N\longrightarrow[0,+\infty)\), con \(g(n)>0\)
per \(n\) sufficientemente grande. Si definiscono
\begin{align*}
  f\in O(g)
  &\iff \exists c>0\ \exists n_0\in\N:\
      f(n)\le c\,g(n)\quad\forall n\ge n_0,\\
  f\in\Omega(g)
  &\iff \exists c>0\ \exists n_0\in\N:\
      f(n)\ge c\,g(n)\quad\forall n\ge n_0,\\
  f\in\Theta(g)
  &\iff \exists c_1,c_2>0\ \exists n_0\in\N:\
      c_1g(n)\le f(n)\le c_2g(n)\quad\forall n\ge n_0.
\end{align*}
\end{definition}

La notazione \(O\) esprime un limite superiore asintotico, \(\Omega\)
un limite inferiore e \(\Theta\) entrambi. In particolare,
\[
  f\in\Theta(g)
  \iff f\in O(g)\ \text{e}\ f\in\Omega(g).
\]
Il significato delle notazioni è distinto dal numero esatto di
operazioni e dalla durata misurata su una macchina.

\begin{example}[Termine dominante]
Per ogni \(n\ge1\),
\[
  3n^2\le 3n^2+10^4n+10^{1000}
  \le (3+10^4+10^{1000})n^2.
\]
Pertanto
\[
  3n^2+10^4n+10^{1000}\in\Theta(n^2).
\]
Analogamente,
\[
  6\in\Theta(1),\qquad
  n^2+3\log_2n+5\in\Theta(n^2),\qquad
  3n^3+10n^2+5n\in\Theta(n^3).
\]
Poiché un'esponenziale di base maggiore di uno cresce più rapidamente
di ogni potenza fissata di \(n\), si ha anche
\[
  3\cdot2^n+5n^{15}+100n\in\Theta(2^n).
\]
\end{example}

\section{Un algoritmo quadratico per l'inversione di una lista}
\label{sec:fi-complessita-reverse-quadratico}

Si consideri il problema di costruire una nuova lista contenente gli
elementi di una lista data in ordine inverso. Gli accessi e le
operazioni sulle liste sono descritti nella \cref{sec:fi-python-liste}.
La seguente implementazione inserisce ogni elemento all'inizio della
lista di uscita, spostando verso destra quelli già inseriti.

\begin{unipdcode}[language=Python]
def reverse_quadratico(lista):
    n = len(lista)
    output = [0] * n
    output_size = 0
    for i in range(n):
        for j in range(output_size, 0, -1):
            output[j] = output[j - 1]
        output[0] = lista[i]
        output_size += 1
    return output
\end{unipdcode}


All'inizio dell'iterazione di indice \(i\), i primi \(i\) elementi di
\texttt{output} contengono, in ordine inverso, i primi \(i\) elementi
della lista di ingresso. Il ciclo interno li sposta di una posizione;
l'assegnamento successivo inserisce \texttt{lista[i]} in posizione
zero. Dopo \(n\) iterazioni si ottiene quindi la lista invertita.

Per determinare il costo si adotta un conteggio semplificato:
l'istruzione \texttt{n = len(lista)} e gli assegnamenti semplici
valgono un passo ciascuno, mentre la costruzione di
\texttt{[0] * n} costa \(n\) passi. Si trascurano i controlli
e gli aggiornamenti impliciti dei cicli e la restituzione finale.
L'inizializzazione costa pertanto \(n+3\) passi.

Nel ciclo interno l'istruzione
\texttt{output[j] = output[j - 1]} comporta una sottrazione,
un accesso in lettura e uno in scrittura: tre passi. Poiché
\texttt{output\_size} vale \(i\) all'inizio dell'iterazione esterna,
questa istruzione viene eseguita \(i\) volte. Le istruzioni successive
costano complessivamente altri tre passi: una lettura, una scrittura e
un incremento. Ne segue
\begin{align*}
  T(n)
    &=n+3+\sum_{i=0}^{n-1}(3i+3)\\
    &=n+3+3\frac{n(n-1)}2+3n\\
    &=\frac32n^2+\frac52n+3
     \in\Theta(n^2).
\end{align*}
L'espressione è esatta nel modello di conteggio adottato; non è il
numero effettivo di istruzioni macchina eseguite da Python. Il costo
quadratico dipende dagli spostamenti ripetuti, non dai valori degli
elementi. La lista di uscita occupa \(\Theta(n)\) spazio.

\section{Un algoritmo lineare per lo stesso problema}
\label{sec:fi-complessita-reverse-lineare}

Gli spostamenti non sono necessari se ciascun elemento viene scritto
direttamente nella propria posizione definitiva. Visitando gli indici
in ordine crescente, si copia l'elemento di indice \(n-1-i\) nella
posizione \(i\) della lista di uscita.

\begin{unipdcode}[language=Python]
def reverse_lineare(lista):
    n = len(lista)
    output = [0] * n
    for i in range(n):
        output[i] = lista[n - 1 - i]
    return output
\end{unipdcode}


L'inizializzazione costa \(n+2\) passi nel modello precedente. Ogni
iterazione esegue due sottrazioni, una lettura e una scrittura: quattro
passi, indipendenti da \(n\). Si ottiene quindi
\[
  T(n)=n+2+4n=5n+2\in\Theta(n).
\]
Entrambe le implementazioni producono lo stesso risultato e richiedono
\(\Theta(n)\) memoria per la lista di uscita, ma la seconda passa da
un costo temporale quadratico a uno lineare.

Per confrontarle sperimentalmente, il seguente programma registra in
\texttt{tempi.csv} la durata delle esecuzioni su liste di dimensione
crescente. Si assumono definite le due funzioni precedenti.

\begin{unipdcode}[language=Python]
import time

with open("tempi.csv", "w", encoding="utf-8") as file:
    file.write("algoritmo,n,secondi\n")
    for n in (500, 1000, 2000, 4000, 8000):
        dati = list(range(n))
        for funzione in (reverse_quadratico, reverse_lineare):
            inizio = time.perf_counter()
            funzione(dati)
            durata = time.perf_counter() - inizio
            file.write(f"{funzione.__name__},{n},{durata:.8f}\n")
\end{unipdcode}


Il file risultante contiene, per ogni misura, il nome dell'algoritmo,
la dimensione dell'ingresso e il tempo in secondi. Le misure
singole restano soggette alle fluttuazioni dell'esecuzione.
Per input sufficientemente grandi, il rapporto teorico fra i costi dei
due algoritmi cresce proporzionalmente a \(n\).

\section{Ordini di crescita e trattabilità}
\label{sec:fi-complessita-ordini}

Fra gli ordini di crescita più comuni si distinguono quello costante
\(\Theta(1)\), logaritmico \(\Theta(\log n)\), lineare \(\Theta(n)\),
lineare-logaritmico \(\Theta(n\log n)\), quadratico \(\Theta(n^2)\)
e, più in generale, polinomiale \(\Theta(n^k)\) per \(k>0\) fissato.
Una crescita esponenziale come \(\Theta(2^n)\) supera asintoticamente
qualsiasi polinomio di grado fissato; la funzione \(n^n\) cresce
ancora più rapidamente.

La distinzione fra crescita polinomiale ed esponenziale è importante
per la trattabilità dei problemi, ma non determina da sola la
praticabilità di un'implementazione: incidono anche la dimensione
concreta dell'ingresso, i coefficienti e le risorse disponibili. In
particolare, un algoritmo polinomiale con esponente elevato può essere
impraticabile per dimensioni moderate. Le applicazioni alla ricerca e all'ordinamento sono
esaminate nel \cref{chap:ricerca-ordinamento}.

\section{Complessità della ricorsione di Fibonacci}
\label{sec:fi-complessita-fibonacci}

Le due implementazioni di Fibonacci sono presentate nella
\cref{sec:fi-fibonacci}, dove si è stabilito che il numero complessivo
delle invocazioni ricorsive è \(C_n=2F_n-1\).
Nel modello a costo unitario, ogni invocazione compie un numero
costante di operazioni, escluse quelle delle chiamate interne. Il
tempo dell'algoritmo ricorsivo è dunque
\[
  T_{\mathrm{ric}}(n)\in\Theta(C_n)=\Theta(F_n).
\]
Dall'\cref{eq:fi-fibonacci} segue \(F_{n+2}\ge2F_n\). Per \(n\ge2\)
si ottengono pertanto le stime
\[
  2^{\lfloor(n-1)/2\rfloor}\le F_n\le2^{n-2},
\]
che mostrano la crescita esponenziale del numero delle invocazioni.
Il costo dipende dal ricalcolo degli stessi sottoproblemi, visibile
nella \cref{fig:fi-fibonacci-albero}.

La versione iterativa, invece, calcola ogni termine una sola volta:
\(T_{\mathrm{iter}}(n)\in\Theta(n)\) e impiega \(\Theta(1)\) spazio
ausiliario. La ricorsione raggiunge profondità \(\Theta(n)\) e richiede
\(\Theta(n)\) spazio per la pila delle chiamate. Come nel confronto
del fattoriale della \cref{sec:fi-pila-ricorsione}, queste stime
contano le operazioni e le variabili senza includere il costo crescente
dell'aritmetica su interi arbitrariamente grandi.
